
\documentclass[12pt,a4paper]{article}

\usepackage[a4paper,margin=1in]{geometry}
\usepackage{graphicx}
\usepackage{booktabs}
\usepackage{amsmath,amssymb}
\usepackage{float}
\usepackage{hyperref}
\usepackage{xcolor}
\usepackage{listings}
\usepackage{enumitem}
\usepackage{fancyhdr}
\usepackage{longtable}

\hypersetup{colorlinks=true,linkcolor=blue,urlcolor=blue}

\pagestyle{fancy}
\fancyhf{}
\lhead{ICS1512}
\rhead{Machine Learning Algorithms Laboratory}
\cfoot{\thepage}

\lstset{
language=Python,
basicstyle=\ttfamily\small,
numbers=left,
frame=single,
breaklines=true,
showstringspaces=false
}

\begin{document}

\begin{tabular}{|p{4cm}|p{11cm}|}
\hline
Experiment No. & 5 \\ \hline
Experiment Title & Decision Tree and Random Forest: A Comparative Classification Study\footnotemark \\ \hline
Student Name & Sanjay S \\ \hline
Register Number & 3122247001056 \\ \hline
Faculty & Dr. Poreddy Ajay Kumar Reddy \\ \hline
Submission Date & 25/08/2026 \\ \hline
\end{tabular}
\footnotetext{\textbf{GitHub Repository: }\url{https://github.com/sanjay-6727/Introdution-to-Machine-Learning-}}

\section{Objective}
\begin{itemize}
\item Implement a Decision Tree classifier from scratch using scikit-learn.
\item Extend the Decision Tree into a Random Forest ensemble and see what bagging actually buys you.
\item Study how depth, splitting, and leaf-size hyperparameters push a tree toward overfitting or
generalization.
\item Pick hyperparameters properly, using 5-fold cross-validation rather than eyeballing the test set.
\item Compare a single tree against the forest on the same held-out data and figure out where the
extra complexity is actually earning its keep.
\end{itemize}

\section{Problem Statement}
Given the Wisconsin Diagnostic Breast Cancer dataset, 30 numeric features computed from digitized
images of breast mass biopsies, classify each sample as malignant or benign. A single Decision Tree
and a Random Forest are both tuned with 5-fold cross-validation and compared on the same 80/20
train-test split, to check how much the ensemble's variance reduction is actually worth on a dataset
this size.

\section{Dataset Description}
\begin{longtable}{|p{6cm}|p{8cm}|}
\hline
Dataset Name & Wisconsin Diagnostic Breast Cancer (WDBC) \\ \hline
Dataset Source & UCI Machine Learning Repository, \url{https://archive.ics.uci.edu} \\ \hline
Number of Samples & 569 \\ \hline
Number of Features & 30 numerical attributes \\ \hline
Number of Classes & 2 (Malignant = M, Benign = B) \\ \hline
Class Balance & 357 Benign (62.74\%), 212 Malignant (37.26\%) \\ \hline
Missing Values & none \\ \hline
Train-Test Split & 80/20 (455 train / 114 test), with 5-fold cross validation for tuning \\ \hline
\end{longtable}

\section{Exploratory Data Analysis}
\begin{itemize}
\item Class distribution is moderately imbalanced, roughly 63:37 in favor of benign, not severe
enough to need resampling but enough to keep an eye on recall for the malignant class.
\item Feature correlation heatmap across all 30 attributes shows the ``mean'', ``error'', and
``worst'' versions of the same measurement (radius, perimeter, area) are, unsurprisingly, very
strongly correlated with each other, since they are just different summary statistics of the same
underlying cell nuclei measurements.
\item Pairplot on the most predictive features shows benign and malignant samples separating into
fairly distinct, if slightly overlapping, clusters, mainly driven by the worst-value and
concave-points features.
\item Histograms, boxplots, and violin plots (generated separately for all 30 features) show most
features are right-skewed with a long tail of outliers on the malignant side, exactly the kind of
distribution a tree-based model handles comfortably since it splits on thresholds rather than
assuming any particular shape.
\item No missing values anywhere in the dataset, confirmed both numerically and visually.
\end{itemize}

\begin{figure}[H]
\centering
\includegraphics[width=0.95\linewidth]{outputs/combined/eda_grid.png}
\caption{Correlation heatmap and pairplot of the most predictive features (top row), outlier counts
per feature and missing-value check across all 30 features (bottom row).}
\end{figure}

\textbf{Inference:}

The heatmap makes it pretty obvious the 30 features are not 30 independent pieces of information,
groups of them are just repackaged versions of the same three or four underlying measurements
(size, concavity, texture). This matters for the Random Forest more than the tree, since
\texttt{max\_features} decides how much of that redundancy each individual tree in the forest gets to
see. The pairplot backs up the correlation story: the classes are separable but not by a single
straight cut, which is exactly the setting where a tree's axis-aligned, piecewise splits should do
reasonably well without needing anything fancier like an SVM kernel.

\section{Data Preprocessing}
\begin{itemize}
\item Diagnosis labels encoded as M = 1 (malignant), B = 0 (benign).
\item No missing values, no categorical columns, so no imputation or encoding of features was
needed.
\item Decision Trees and Random Forests split on raw thresholds, so unlike Experiment 4 there was no
need for \texttt{StandardScaler}, feature scale simply does not affect where a tree decides to split.
\item 80/20 stratified train-test split, with the tuning itself done through 5-fold cross-validation
on the training set only, so the test set stays untouched until final evaluation.
\end{itemize}

\begin{lstlisting}[caption={5-fold cross-validated grid search for the Decision Tree and Random Forest}]
dt_param_grid = {
    "criterion": ["gini", "entropy"],
    "max_depth": [2, 3, 4, 5, 6, 8, 10, None],
    "min_samples_split": [2, 5, 10],
    "min_samples_leaf": [1, 2, 4],
}

rf_param_grid = {
    "n_estimators": [50, 100, 200],
    "max_depth": [3, 5, 8, 10, None],
    "max_features": ["sqrt", "log2"],
    "bootstrap": [True, False],
}

for params in ParameterGrid(dt_param_grid):
    clf = DecisionTreeClassifier(random_state=42, **params)
    scores = cross_val_score(clf, X_train, y_train, cv=5, scoring="accuracy")
    # record mean/std CV accuracy per combination

for params in ParameterGrid(rf_param_grid):
    clf = RandomForestClassifier(random_state=42, n_jobs=-1, **params)
    scores = cross_val_score(clf, X_train, y_train, cv=5, scoring="accuracy")
    # record mean/std CV accuracy per combination
\end{lstlisting}

\section{Machine Learning Algorithm}

\textbf{Decision Tree Classifier.} A Decision Tree recursively partitions the feature space by
picking, at each node, the feature and threshold that gives the largest reduction in impurity, either
Gini index $G = 1 - \sum_i p_i^2$ or entropy $H = -\sum_i p_i \log_2 p_i$. Splitting continues until a
stopping condition is hit (max depth, minimum samples per split or leaf), producing an interpretable
set of if-else rules. Left unconstrained, a tree will keep splitting until every leaf is pure, which
usually means memorizing noise along with signal, hence the well-known tendency toward high variance
and overfitting.

\textbf{Random Forest Classifier.} A Random Forest builds many such trees, each on a bootstrapped
sample of the training data (bagging), and at every split each tree is only allowed to consider a
random subset of the features (controlled by \texttt{max\_features}). Predictions are combined by
majority vote across all trees. Averaging over many trees that are individually high-variance but
decorrelated from each other tends to cancel out a lot of that variance, which is the whole point of
the ensemble: it trades a bit of per-tree interpretability for a model that generalizes more
reliably.

\section{Experimental Setup}
\begin{tabular}{ll}
\toprule
Python Version & 3.11.9 \\
Scikit-Learn Version & 1.7.2 \\
IDE & Anaconda Navigator / VS Code \\
Operating System & Windows 11 \\
Hardware & \\
\bottomrule
\end{tabular}

\section{Hyperparameter Selection}
\begin{tabular}{p{2.3cm}p{12.5cm}}
\toprule
Parameter & Value\\
\midrule
DT grid & criterion in \{gini, entropy\}, max\_depth in \{2,3,4,5,6,8,10,None\}, min\_samples\_split
in \{2,5,10\}, min\_samples\_leaf in \{1,2,4\} \\
RF grid & n\_estimators in \{50,100,200\}, max\_depth in \{3,5,8,10,None\}, max\_features in
\{sqrt,log2\}, bootstrap in \{True,False\} \\
\bottomrule
\end{tabular}

\vspace{0.3cm}
\begin{tabular}{p{3cm}p{2.6cm}p{7cm}p{2cm}}
\toprule
Model & Search Method & Best Parameters & Best CV Acc. \\
\midrule
Decision Tree & 5-Fold CV Grid Search & criterion=entropy, max\_depth=4, min\_samples\_split=2,
min\_samples\_leaf=2 & 94.07\% \\
Random Forest & 5-Fold CV Grid Search & n\_estimators=200, max\_depth=None, max\_features=sqrt,
bootstrap=False & 97.14\% \\
\bottomrule
\end{tabular}

\subsection{Decision Tree Cross-Fold Results (Top Combinations)}
\begin{longtable}{p{2.3cm}p{2.3cm}p{3.3cm}p{3cm}p{3.3cm}}
\toprule
Criterion & Max Depth & Avg CV Acc. (\%) & Std CV Acc. & Avg Train Acc. (\%) \\
\midrule
\endhead
entropy & 4 & 94.07 & 2.15 & 98.19 \\
entropy & 4 & 94.07 & 2.15 & 98.19 \\
entropy & 4 & 93.85 & 1.79 & 97.86 \\
entropy & 4 & 93.85 & 1.64 & 98.46 \\
entropy & 4 & 93.63 & 2.01 & 98.30 \\
entropy & 5 & 93.63 & 2.45 & 98.30 \\
gini & 6 & 93.19 & 3.06 & 98.74 \\
\bottomrule
\end{longtable}

\subsection{Random Forest Cross-Fold Results (Top Combinations)}
\begin{longtable}{p{2.3cm}p{2.3cm}p{2.6cm}p{2.3cm}p{3.3cm}}
\toprule
n\_estimators & Max Depth & Max Features & Bootstrap & Avg CV Accuracy (\%) \\
\midrule
\endhead
200 & 5 & sqrt & False & 97.14 \\
50 & 10 & sqrt & False & 97.14 \\
200 & 10 & sqrt & False & 97.14 \\
200 & None & sqrt & False & 97.14 \\
100 & 10 & log2 & False & 96.92 \\
100 & 5 & sqrt & False & 96.92 \\
50 & None & log2 & False & 96.92 \\
\bottomrule
\end{longtable}

A quick pattern worth noting: for the Random Forest, \texttt{bootstrap=False} (each tree sees the
full training set, only the feature subsampling stays random) consistently outperformed
\texttt{bootstrap=True} across almost every other setting held fixed. With only 455 training rows,
bootstrapping away roughly a third of them per tree seems to hurt more than the extra
decorrelation helps, whereas on a much larger dataset the usual bagging advantage would likely show
up again.

\section{Results}

\begin{tabular}{lc}
\toprule
Decision Tree Metric & Value \\
\midrule
Accuracy & 0.9298 \\
Precision & 1.0000 \\
Recall & 0.8095 \\
F1 Score & 0.8947 \\
ROC-AUC & 0.9563 \\
\bottomrule
\end{tabular}
\hspace{1cm}
\begin{tabular}{lc}
\toprule
Random Forest Metric & Value \\
\midrule
Accuracy & 0.9737 \\
Precision & 1.0000 \\
Recall & 0.9286 \\
F1 Score & 0.9630 \\
ROC-AUC & 0.9940 \\
\bottomrule
\end{tabular}

\vspace{0.3cm}
\textbf{5-Fold Cross-Validation Accuracy Comparison}

\begin{tabular}{lcccccc}
\toprule
Model & Fold 1 & Fold 2 & Fold 3 & Fold 4 & Fold 5 & Average \\
\midrule
Decision Tree & 0.9341 & 0.9560 & 0.9011 & 0.9560 & 0.9560 & 0.9407 \\
Random Forest & 0.9670 & 0.9780 & 0.9560 & 0.9670 & 0.9890 & 0.9714 \\
\bottomrule
\end{tabular}

\vspace{0.3cm}
\textbf{Comparative Analysis}

\begin{tabular}{lcc}
\toprule
Criterion & Decision Tree & Random Forest \\
\midrule
Test Accuracy & 0.9298 & 0.9737 \\
CV Accuracy (mean) & 0.9407 & 0.9714 \\
Model Complexity & Low (single tree) & High (200 trees) \\
Training Time & Very Low & Moderate \\
Interpretability & High & Low \\
\bottomrule
\end{tabular}

\vspace{0.3cm}
\textbf{Top Random Forest Feature Importances}

\begin{tabular}{lc}
\toprule
Feature & Importance \\
\midrule
worst perimeter & 0.1536 \\
worst area & 0.1287 \\
worst concave points & 0.1050 \\
worst radius & 0.0883 \\
mean concave points & 0.0846 \\
\bottomrule
\end{tabular}

\section{Result Visualization}

\begin{figure}[H]
\centering
\includegraphics[width=0.9\linewidth]{outputs/combined/model_grid.png}
\caption{Trained Decision Tree structure, criterion=entropy, max\_depth=4 (top), with test-set
confusion matrices for the Decision Tree (bottom left) and Random Forest (bottom right).}
\end{figure}

\begin{figure}[H]
\centering
\includegraphics[width=0.9\linewidth]{outputs/combined/results_grid.png}
\caption{ROC curves and Random Forest feature importances (top row), with the full
Decision Tree vs. Random Forest metric comparison (bottom).}
\end{figure}

\textbf{Inference:}

Both models get precision-1.0 on the test set, meaning neither one ever calls a benign case
malignant, which is the easy direction to get right given the 63:37 class split. The real
difference shows up in recall on the malignant class: the single tree misses about 1 in 5 malignant
cases (recall 0.8095) while the forest catches nearly all of them (recall 0.9286). That is visible
directly in the confusion matrices too, the Decision Tree panel shows a noticeably thicker
false-negative column than the Random Forest one. The ROC curves tell the same story from a
threshold-independent angle, Random Forest's curve hugs the top-left corner much more tightly
(AUC 0.9940 against the tree's 0.9563). The feature importance chart lines up neatly with the EDA,
the worst-value and concave-points features that stood out in the pairplot are exactly the ones the
forest leans on most.

\section{Discussion}
Random Forest wins on every single metric here, higher CV accuracy (97.14\% vs 94.07\%), higher test
accuracy (0.9737 vs 0.9298), and a big jump in recall and F1 for the malignant class. None of that is
surprising given the theory: a single tree tuned to max\_depth=4 is already a fairly shallow,
low-variance tree by Decision Tree standards, and even so, averaging 200 of them still buys a
meaningful drop in test error. What is somewhat more interesting is that the best Random Forest
configuration used \texttt{bootstrap=False}, so essentially all of the improvement over the single
tree is coming from the random feature subsampling at each split (decorrelating the trees) rather
than from bagging the training rows, which on a dataset this size (455 training rows) does not have
much room to help. For a medical screening context specifically, the gap in recall matters more than
the accuracy gap on its own: missing a malignant tumour is a far more expensive mistake than a false
alarm, and that is exactly the axis on which Random Forest pulls ahead.

\section{Observation Questions}
\begin{itemize}
\item \textbf{How does tree depth affect overfitting in Decision Trees?} In the cross-fold results,
average train accuracy keeps climbing toward 100\% as max\_depth increases past 6 or 8, while average
CV accuracy actually peaks around max\_depth=4 and then drifts slightly downward. That gap between a
rising train score and a flat-to-falling CV score is overfitting playing out directly in the numbers,
a deeper tree keeps carving the training data into smaller and smaller pure leaves that do not
generalize any further.
\item \textbf{Which hyperparameter had the greatest impact on performance?} For the Decision Tree, it
was clearly max\_depth, criterion (gini vs entropy) and the leaf/split minimums only shifted CV
accuracy by fractions of a percent once depth was reasonable. For the Random Forest,
\texttt{bootstrap} had the most consistent effect (False beating True almost everywhere in the grid),
with \texttt{n\_estimators} mattering a bit less once it got past 50 trees.
\item \textbf{How does Random Forest improve generalization?} By training many trees on
different bootstrapped samples and random feature subsets and then averaging their votes, the
forest cancels out a lot of the idiosyncratic errors any single tree would make on its own training
split. Individual trees can still overfit their own bootstrap sample, but their mistakes tend not to
line up with each other, so the majority vote ends up closer to the true decision boundary than any
one tree.
\item \textbf{Did ensemble learning always improve performance? Why or why not?} Here, yes, across
every metric and every fold. But it is not automatic in general: ensembling helps most when the base
trees are reasonably accurate individually and their errors are not too correlated. On a dataset
that is small, extremely noisy, or where the trees are so shallow they underfit badly, or so
correlated that bagging brings no diversity, a Random Forest can end up barely better than a single
well-tuned tree, or occasionally worse if it is tuned carelessly.
\end{itemize}

\section{Conclusion}
Both a Decision Tree and a Random Forest were implemented and tuned on the WDBC dataset using 5-fold
cross-validation, with hyperparameters chosen on cross-validation accuracy rather than the test set.
The Random Forest came out ahead on every metric (97.37\% test accuracy and 0.9940 ROC-AUC against the
tree's 92.98\% and 0.9563), confirming the textbook claim that bagging plus random feature selection
reduces variance and improves generalization compared to a single tree. The single Decision Tree
remains useful mainly for its interpretability, the whole decision logic fits in one diagram, while
the forest is the model to actually deploy if recall on the malignant class is the priority.

\section{Learning Outcomes}
\begin{itemize}
\item Saw overfitting show up concretely in the cross-fold table, not just as a definition, rising
train accuracy against a flattening CV accuracy as max\_depth increased.
\item Learned that for a Random Forest, the two randomness sources (bootstrapping rows vs.
subsampling features per split) do not always contribute equally, on this dataset the feature
subsampling was doing most of the work.
\item Got a clearer sense of why accuracy alone can be misleading on an imbalanced dataset, recall on
the minority (malignant) class told a much sharper story than the overall accuracy numbers did.
\item Practiced picking hyperparameters purely from cross-validation performance and only touching
the test set once, at the very end, for a fair final comparison.
\end{itemize}

\section{References}
\begin{enumerate}
\item Scikit-learn: Decision Trees, \url{https://scikit-learn.org/stable/modules/tree.html}.
\item Scikit-learn: Random Forest, \url{https://scikit-learn.org/stable/modules/ensemble.html\#random-forests}.
\item UCI Machine Learning Repository, Wisconsin Diagnostic Breast Cancer Dataset,
\url{https://archive.ics.uci.edu}.
\item Stephen Marsland, ``Machine Learning - An Algorithmic Perspective'', 2nd Edition, Chapman and
Hall/CRC Machine Learning and Pattern Recognition Series, 2015.
\item Ethem Alpaydin, ``Introduction to Machine Learning'', 3rd Edition, The MIT Press, 2014.
\end{enumerate}

\end{document}
